\documentclass[conference]{IEEEtran}
\IEEEoverridecommandlockouts

\usepackage{cite}
\usepackage{amsmath,amssymb,amsfonts}
\usepackage{graphicx}
\usepackage{booktabs}
\usepackage{multirow}
\usepackage{url}
\usepackage{xcolor}
\usepackage[hidelinks]{hyperref}

% Placeholder marker for values the authors must fill before submission.
\newcommand{\tofill}[1]{\textcolor{red}{[FILL: #1]}}

% Rupee symbol without external font dependencies. Swap for \usepackage{tfrupee}
% and its \rupee glyph if the venue's toolchain supports it.
\newcommand{\rupee}{Rs.\,}

\begin{document}

\title{TrustNote-AI: A Frozen-Backbone MobileNetV3 Classifier for\\
Counterfeit Indian Banknote Screening on Commodity Hardware}

\author{
\IEEEauthorblockN{Sai Shiva Ganesh}
\IEEEauthorblockA{\textit{Department of Computer Science and Engineering}\\
\textit{SRM Institute of Science and Technology}\\
Chennai, India\\
sg3142@srmist.edu.in}
\and
\IEEEauthorblockN{Sricharan Suresh}
\IEEEauthorblockA{\textit{Department of Computer Science and Engineering}\\
\textit{SRM Institute of Science and Technology}\\
Chennai, India\\
ss1833@srmist.edu.in}
}

\maketitle

\begin{abstract}
Counterfeit banknote screening in low-resource settings runs on phone cameras
and shared desktops, not on ultraviolet scanners. We present TrustNote-AI, a
binary REAL/FAKE classifier for Indian banknotes built on MobileNetV3-Large with
ImageNet-pretrained weights. The convolutional backbone stays frozen throughout
training; only a 0.56~M-parameter classification head learns. On a held-out test
set of 1121 images the model reaches 96.97\% accuracy, 0.9555 $F_1$ on the
counterfeit class, and 0.9959 ROC-AUC, while occupying 14.3~MB in 32-bit
floating point and classifying an image in 9.1~ms on a desktop CPU. No GPU takes
part at any stage: the reported run trained for nine epochs in 36.09 minutes on
an Intel Core i5-1335U laptop CPU before early stopping, with the best
validation $F_1$ of 0.9775 recorded at epoch four. The result carries a caveat
that we treat as a central finding rather than a footnote: a frozen ImageNet
feature extractor,
which has never seen a security thread or an intaglio print, separates the two
classes at 0.9959 AUC. Coarse photometric and print-process cues therefore drive
the decision, not the security features a human examiner reads. We report the
confusion matrix in full and specify the acquisition-controlled evaluation
protocol that any deployment claim would require.
\end{abstract}

\begin{IEEEkeywords}
counterfeit detection, Indian currency, transfer learning, MobileNetV3,
edge inference, image classification, dataset bias
\end{IEEEkeywords}

\section{Introduction}

The Reserve Bank of India reports counterfeit notes detected across the banking
system every year, and the detected figure understates circulation because notes
that never reach a bank counter never enter the count~\cite{rbi_annual_report}.
Verification hardware closes part of the gap. A ultraviolet lamp or a
magnetic-ink reader costs more than a small merchant will spend and travels
poorly, so the transactions most exposed to counterfeits are screened by eye or
not at all.

A classifier that runs on a phone changes the economics of that screening. This
paper reports what such a classifier achieves, and what its headline accuracy
conceals.

We train MobileNetV3-Large~\cite{howard2019mobilenetv3} on a
denomination-stratified corpus of real and counterfeit Indian banknote
photographs and evaluate it on a held-out 15\% partition. Three design decisions
define the system. The backbone remains frozen, so training touches 16\% of the
parameters. A weighted sampler corrects a 2:1 class imbalance during training
while the test partition retains its natural prior. The augmentation pipeline
models photographic nuisance variables, including perspective, blur, and
occlusion, rather than generic geometric jitter.

The contributions of this paper are:

\begin{enumerate}
\item A complete, reproducible pipeline for binary Indian banknote
authentication that trains in under ten epochs and deploys in 14.3~MB, with
audit, split, training, evaluation, and inference stages released as code.
\item A measured result on 1121 held-out images: 96.97\% accuracy and 0.9959
ROC-AUC, reported with the full confusion matrix and per-class breakdown rather
than accuracy alone.
\item An analysis showing that a frozen ImageNet backbone suffices for this
performance, and an argument that this fact bounds what the reported numbers can
be claimed to demonstrate about counterfeit detection as opposed to image-source
discrimination.
\item An evaluation protocol, specified in Section~\ref{sec:protocol}, that
would separate the two.
\end{enumerate}

\section{Related Work}

Automated banknote authentication divides into three families.

\textbf{Sensor-based methods} read physical properties: ultraviolet
fluorescence, infrared absorption, magnetic ink response, and paper thickness.
These methods target the substrate that counterfeiters find hardest to
reproduce, and they achieve the accuracy that central banks require. They also
require dedicated hardware, which rules them out for the deployment target
considered here.

\textbf{Hand-engineered visual features} extract security elements from
photographs: the security thread, the latent image, microlettering, the
watermark region, and the intaglio texture of the portrait. Systems in this
family localize a region of interest, apply a filter bank or texture descriptor,
and classify the response. Performance depends on localization accuracy, which
degrades under the perspective and lighting variation typical of handheld
capture~\cite{placeholder_indian_currency_survey}.

\textbf{Learned representations} replace the feature engineering with a
convolutional network. Transfer learning from ImageNet~\cite{deng2009imagenet}
dominates this family because banknote corpora are small by computer vision
standards. Yosinski et al.~\cite{yosinski2014transferable} established that
early convolutional layers transfer across tasks with little loss, which
motivates freezing them when target data is scarce. MobileNetV2 and
V3~\cite{sandler2018mobilenetv2,howard2019mobilenetv3} extend the approach to
mobile inference budgets through inverted residual blocks and
squeeze-and-excitation attention~\cite{hu2018squeeze}.

Work in the third family reports accuracies above 95\% with regularity. Torralba
and Efros~\cite{torralba2011bias} give the reason to read those numbers with
care: when the two classes in a corpus were photographed under different
conditions, a classifier separates the conditions and the reported accuracy
measures nothing about the intended task. Section~\ref{sec:discussion} applies
that argument to our own result.

\section{Dataset}

\subsection{Composition}

The corpus contains photographs of Indian banknotes in seven denominations
(\rupee10, \rupee20, \rupee50, \rupee100, \rupee200, \rupee500, and \rupee2000),
organized as a two-level hierarchy of class and denomination. The prediction
task collapses the denomination axis: the label space is
$\mathcal{Y} = \{\texttt{REAL}, \texttt{FAKE}\}$, encoded as
$\texttt{REAL} \mapsto 0$ and $\texttt{FAKE} \mapsto 1$, with \texttt{FAKE}
treated as the positive class for precision and recall.

The corpus holds 7473 images, 4960 labeled \texttt{REAL} and 2513 labeled
\texttt{FAKE}. A stratified 70/15/15 split produces the partitions in
Table~\ref{tab:splits}. The \texttt{REAL}-to-\texttt{FAKE} ratio holds at
1.97:1 across all three.

\begin{table}[t]
\caption{Partition sizes.}
\label{tab:splits}
\centering
\small
\begin{tabular}{@{}lrrr@{}}
\toprule
\textbf{Partition} & \textbf{\texttt{REAL}} & \textbf{\texttt{FAKE}} &
\textbf{Total} \\
\midrule
Training   & 3{,}472 & 1{,}759 & 5{,}231 \\
Validation & 744     & 377     & 1{,}121 \\
Testing    & 744     & 377     & 1{,}121 \\
\midrule
\textbf{Total} & \textbf{4{,}960} & \textbf{2{,}513} & \textbf{7{,}473} \\
\bottomrule
\end{tabular}
\end{table}

The test partition stayed out of both training and model selection. Every number
in Section~\ref{sec:results} comes from a single pass over it after the
checkpoint was fixed.

\subsection{Audit}

Every image passes an audit stage before splitting. The audit walks the
directory tree, computes a SHA-256 digest per file, opens each image with a
verification pass, and records width, height, color mode, and container format.
The stage emits three artifacts: a per-image inventory, a list of files that
failed to decode, and a list of digest collisions.

Exact-duplicate detection at the file level catches copied files. It does not
catch two photographs of the same physical note taken seconds apart, which
produce different digests and near-identical content.
Section~\ref{sec:limitations} returns to this.

\subsection{Partitioning}

The split runs per class and per denomination with a fixed seed of 42. Within
each of the fourteen leaf directories, a 70\% training partition is drawn first,
and the 30\% remainder is halved into validation and test. Stratifying at the
denomination level keeps the denomination distribution stable across the three
partitions, which prevents the classifier from being evaluated on a denomination
mix it never trained on.

The split unit is the image file. It is not the physical note, and it is not the
capture session.

\section{Method}

\subsection{Preprocessing}

All images convert to RGB, resize to $224 \times 224$, and normalize with
ImageNet channel statistics ($\mu = [0.485, 0.456, 0.406]$,
$\sigma = [0.229, 0.224, 0.225]$). Matching the normalization used to pretrain
the backbone keeps the frozen features in the input distribution they were
fitted on.

Resizing to a fixed square discards the aspect ratio. Indian denominations
differ in physical dimensions, so this step removes a cue that would otherwise
identify the denomination. The binary REAL/FAKE task does not need that cue, and
withholding it stops the head from routing its decision through denomination.

\subsection{Augmentation}

Training images pass through a pipeline that models handheld capture rather than
generic geometric perturbation. Table~\ref{tab:aug} lists the operations and the
nuisance variable each one targets.

\begin{table}[t]
\caption{Training augmentation pipeline. Validation and test images receive
resize and normalization only.}
\label{tab:aug}
\centering
\small
\begin{tabular}{@{}lll@{}}
\toprule
\textbf{Operation} & \textbf{Parameters} & \textbf{Models} \\
\midrule
Resize            & $256 \times 256$                & --- \\
RandomResizedCrop & scale $[0.85, 1.0]$             & camera distance \\
RandomRotation    & $\pm 5^{\circ}$                 & handheld tilt \\
RandomAffine      & shift 3\%, shear $2^{\circ}$    & framing drift \\
RandomPerspective & distortion 0.15, $p{=}0.30$     & off-axis capture \\
ColorJitter       & brightness/contrast 0.15        & lighting \\
GaussianBlur      & $k{=}3$, $p{=}0.20$             & motion blur \\
RandomErasing     & 2--8\% area, $p{=}0.15$         & finger, shadow \\
\bottomrule
\end{tabular}
\end{table}

The parameter ranges stay narrow. Rotation caps at five degrees and color jitter
at 15\%, because a banknote photographed for verification is held roughly
square-on and roughly in frame. Wider ranges would inject variation the
deployment distribution does not contain. Random
erasing~\cite{zhong2020randomerasing} applies after tensor conversion and before
normalization, so the erased region carries raw-scale noise.

\subsection{Architecture}

The backbone is MobileNetV3-Large~\cite{howard2019mobilenetv3} with the default
ImageNet weights. The original classifier is discarded and replaced with the
head in Table~\ref{tab:head}, which takes the 960-dimensional pooled feature
vector and emits two logits.

\begin{table}[t]
\caption{Replacement classification head. All parameters listed are trainable;
the backbone is frozen.}
\label{tab:head}
\centering
\small
\begin{tabular}{@{}llr@{}}
\toprule
\textbf{Layer} & \textbf{Shape} & \textbf{Parameters} \\
\midrule
Linear      & $960 \to 512$ & 492{,}032 \\
BatchNorm1d & 512           & 1{,}024 \\
Hardswish   & ---           & 0 \\
Dropout     & $p = 0.30$    & 0 \\
Linear      & $512 \to 128$ & 65{,}664 \\
BatchNorm1d & 128           & 256 \\
Hardswish   & ---           & 0 \\
Dropout     & $p = 0.20$    & 0 \\
Linear      & $128 \to 2$   & 258 \\
\midrule
\multicolumn{2}{@{}l}{\textbf{Trainable total}} & \textbf{559{,}234} \\
\multicolumn{2}{@{}l}{Frozen backbone}          & 2{,}971{,}952 \\
\bottomrule
\end{tabular}
\end{table}

Hardswish activations match the backbone. Batch normalization before each
nonlinearity stabilizes the head while the features feeding it stay fixed.
Dropout tapers from 0.30 to 0.20 across the two hidden layers.

The trainable fraction is 559{,}234 of 3{,}531{,}186 parameters, or 15.84\%. We
fitted no convolutional filter on currency data. All of them arrived from
ImageNet and stayed there.

\subsection{Class Imbalance}

The training loader draws through a \texttt{WeightedRandomSampler} with weights
inversely proportional to class frequency, sampling with replacement for one
epoch-equivalent of draws. Each epoch therefore presents a balanced stream, and
the majority \texttt{REAL} class stops dominating the gradient.

Validation and test loaders sample sequentially without weighting. The reported
metrics reflect the natural 2:1 prior rather than a balanced one.

\subsection{Optimization}

Training minimizes cross-entropy with AdamW~\cite{loshchilov2019adamw} at an
initial learning rate of $10^{-4}$ and weight decay of $10^{-4}$, decayed by a
cosine schedule~\cite{loshchilov2017sgdr} toward $\eta_{\min} = 10^{-6}$ over a
30-epoch horizon. Batch size is 32. Gradients clip to unit $\ell_2$ norm. The
trainer wraps the forward pass in an automatic mixed precision
scope~\cite{micikevicius2018mixed} with a gradient scaler that unscales before
the clip, so the threshold applies to true gradient magnitudes whenever that
path is active.

Model selection tracks validation $F_1$ on the \texttt{FAKE} class rather than
accuracy, which would be inflated by the majority class. Training halts when
validation $F_1$ fails to improve for five consecutive epochs, and the
checkpoint with the best validation $F_1$ is restored for testing. We fix random
seeds at 42 for PyTorch and CUDA.

The pipeline gates mixed precision on \texttt{torch.cuda.is\_available()} and
otherwise falls back to full-precision CPU execution. The reported run took the
CPU path on an Intel Core i5-1335U and finished nine epochs in 36.09 minutes, so
the mixed-precision branch stayed inactive and the scaler ran as a no-op.
Inference carries no CUDA dependency either, and the latency figures in
Section~\ref{sec:cost} come from a CPU-only build.

Training the full pipeline in 36 minutes on a mid-range laptop CPU follows from
the frozen backbone. Gradients flow through 559{,}234 parameters rather than
3.5~M, which puts reproduction of this work within reach of any machine that can
run the dataset through a forward pass.

\section{Results}
\label{sec:results}

\subsection{Convergence}

Training stopped at epoch nine of a 30-epoch budget. Table~\ref{tab:history}
gives the trajectory. The best validation $F_1$ of 0.9775 arrived at epoch four;
the five epochs that followed produced no improvement, and early stopping
triggered.

\begin{table}[t]
\caption{Training history. The selected checkpoint is epoch 4.}
\label{tab:history}
\centering
\small
\begin{tabular}{@{}rrrrr@{}}
\toprule
\textbf{Epoch} & \textbf{Train loss} & \textbf{Val loss} &
\textbf{Val acc.} & \textbf{Val $F_1$} \\
\midrule
1 & 0.2528 & 0.1368 & 0.9436 & 0.9212 \\
2 & 0.1506 & 0.0900 & 0.9687 & 0.9550 \\
3 & 0.1062 & 0.0757 & 0.9758 & 0.9649 \\
\textbf{4} & \textbf{0.0864} & \textbf{0.0530} & \textbf{0.9848} &
\textbf{0.9775} \\
5 & 0.0947 & 0.0529 & 0.9821 & 0.9738 \\
6 & 0.0758 & 0.0466 & 0.9839 & 0.9763 \\
7 & 0.0770 & 0.0533 & 0.9830 & 0.9750 \\
8 & 0.0638 & 0.0464 & 0.9794 & 0.9700 \\
9 & 0.0598 & 0.0453 & 0.9830 & 0.9749 \\
\bottomrule
\end{tabular}
\end{table}

Two features of the trajectory deserve comment. Validation $F_1$ passes 0.92 in
a single epoch, before the head has seen the training set twice. And after epoch
four, training loss keeps falling while validation $F_1$ plateaus, the standard
signature of a head beginning to fit noise that the frozen features cannot
resolve.

\subsection{Test Performance}

Table~\ref{tab:metrics} reports threshold-dependent metrics at the default
argmax decision rule, equivalent to a 0.5 probability threshold.

\begin{table}[t]
\caption{Test performance on 1121 held-out images. Precision, recall, and $F_1$
are computed with \texttt{FAKE} as the positive class.}
\label{tab:metrics}
\centering
\small
\begin{tabular}{@{}lr@{}}
\toprule
\textbf{Metric} & \textbf{Value} \\
\midrule
Accuracy   & 0.9697 \\
Precision  & 0.9432 \\
Recall     & 0.9682 \\
$F_1$      & 0.9555 \\
ROC-AUC    & 0.9959 \\
\bottomrule
\end{tabular}
\end{table}

\begin{table}[t]
\caption{Confusion matrix, 1121 test images.}
\label{tab:cm}
\centering
\small
\begin{tabular}{@{}llrr@{}}
\toprule
& & \multicolumn{2}{c}{\textbf{Predicted}} \\
\cmidrule(l){3-4}
& & \textbf{REAL} & \textbf{FAKE} \\
\midrule
\multirow{2}{*}{\textbf{Actual}}
& \textbf{REAL} & 722 & 22 \\
& \textbf{FAKE} & 12  & 365 \\
\bottomrule
\end{tabular}
\end{table}

The error budget is 34 images out of 1121. Twenty-two genuine notes were flagged
as counterfeit, and twelve counterfeits passed as genuine.

The asymmetry matters for deployment. A false positive sends a genuine note to
manual inspection, which costs a merchant time. A false negative accepts a
counterfeit, which costs the merchant the face value of the note. The current
threshold produces almost twice as many false positives as false negatives, and
the ROC-AUC of 0.9959 indicates headroom to move the operating point further in
that direction. Section~\ref{sec:protocol} treats threshold selection as an
open item rather than a solved one.

\subsection{Deployment Cost}
\label{sec:cost}

The trained network holds 3{,}531{,}186 parameters, of which the replacement head
contributes 559{,}234 (15.84\%). The serialized checkpoint occupies 14.34~MB on
disk, against 14.12~MB of raw 32-bit parameter data.

Table~\ref{tab:latency} reports single-image CPU latency measured on the
released checkpoint. We loaded the weights with strict state-dict matching,
warmed up for ten forward passes, and timed fifty passes at batch size one.

\begin{table}[t]
\caption{Single-image inference cost, CPU only, batch size 1, $224 \times 224$
input, PyTorch 2.10 on an AMD Ryzen 9 7900X with 12 threads.}
\label{tab:latency}
\centering
\small
\begin{tabular}{@{}lr@{}}
\toprule
\textbf{Quantity} & \textbf{Value} \\
\midrule
Parameters (total)   & 3{,}531{,}186 \\
Parameters (head)    & 559{,}234 \\
Checkpoint size      & 14.34~MB \\
Latency, median      & 9.1~ms \\
Latency, 95th pct.   & 17.6~ms \\
Latency, min / max   & 7.9 / 64.6~ms \\
\bottomrule
\end{tabular}
\end{table}

These figures come from a desktop part, while the training run used a laptop
i5-1335U. A 9.1~ms median leaves room for the gap: even a tenfold penalty on a
phone SoC keeps a single verification under 100~ms. INT8 post-training
quantization~\cite{jacob2018quantization} would cut the footprint to roughly
3.6~MB, and we have measured neither its latency nor its effect on accuracy.

\section{Discussion}
\label{sec:discussion}

\subsection{What the Frozen Backbone Implies}

The central result of this paper is that 96.97\% accuracy arrived without
updating a single convolutional filter.

MobileNetV3's ImageNet features encode edges, textures, color distributions, and
object-level shape statistics. Nothing in the pretraining corpus contains a
security thread, an intaglio ridge, microlettering, or a watermark. A trained
human examiner separates genuine notes from counterfeits on those four cues, and
they demand the fine-grained, spatially localized evidence that a frozen generic
feature extractor is least equipped to supply.

A linear-plus-MLP probe on those frozen features reaches 0.9959 ROC-AUC. The
straightforward reading is that the two classes in this corpus are separable by
coarse, global, low-level statistics. Candidate sources for such a signal
include print process (offset or inkjet reproduction versus intaglio), paper
reflectance, color gamut differences from the reproduction pipeline, and the
capture conditions under which each class was photographed.

Some of those are legitimate counterfeit evidence. Print-process artifacts are
real forensic signal, and a model keying on them is doing useful work. Capture
conditions are not. Should the two classes have been photographed on different
equipment or under different room lighting, the classifier reaches these numbers
by identifying the photograph rather than the note~\cite{torralba2011bias}, and
the reported metrics transfer to no deployment at all.

Our evidence does not separate the two cases, so we describe the system as a
class separator for this corpus and stop short of calling it a counterfeit
detector.

\subsection{Diagnostic Experiments}
\label{sec:protocol}

Four experiments would resolve the question, ordered by cost.

\textbf{Grad-CAM attribution.} Compute class activation
maps~\cite{selvaraju2017gradcam} over the test partition and measure what
fraction of the attribution mass falls on security-feature regions (thread,
watermark window, portrait intaglio) against background and margin. Attribution
concentrated on borders or uniform paper indicates capture-artifact reliance.
This costs an afternoon and should run first.

\textbf{Backbone fine-tuning.} Unfreeze the backbone and fine-tune at a reduced
learning rate. The code already exposes \texttt{unfreeze\_backbone()}, and the
comment in \texttt{train.py} labels the current run ``Stage 1,'' so the two-stage
schedule was designed and left unexecuted. If fine-tuning yields little gain,
the frozen features already saturate the available signal, which supports the
coarse-cue hypothesis.

\textbf{Cross-source evaluation.} Partition by capture session or device rather
than by file, train on one source, and test on a held-out source. A drop from
0.97 accuracy toward chance under this protocol identifies source
discrimination. This is the decisive experiment, and it requires acquisition
metadata the current corpus does not carry.

\textbf{Note-level splitting.} Group all photographs of each physical note into
a single partition. File-level splitting places two shots of one note into
training and test, which inflates the reported metrics by an amount we cannot
currently bound.

\subsection{Threshold Selection}

The system decides at the argmax, which is the correct rule only when the two
error types carry equal cost and the softmax outputs are calibrated. Neither
holds. A false negative costs the note's face value; a false positive costs a
few seconds. Modern networks also report overconfident
softmax scores~\cite{guo2017calibration}, so a merchant should not read the 0.99
confidence in our example output as a 1-in-100 error rate without temperature
scaling fitted on the validation partition.

A deployed version should fit a temperature on validation data, then select the
threshold that minimizes expected cost under a stated cost ratio, and report
performance at that operating point rather than at 0.5.

\section{Limitations}
\label{sec:limitations}

\begin{itemize}
\item \textbf{Provenance.} The corpus lacks acquisition metadata, so we cannot
verify that genuine and counterfeit images share capture conditions. This is the
binding limitation on every number in Section~V.

\item \textbf{Splitting granularity.} Partitioning by file rather than by
physical note permits near-duplicate leakage between training and test. The
audit stage detects byte-identical files only.

\item \textbf{Counterfeit quality.} The corpus does not stratify counterfeits by
sophistication. Performance against high-quality reproductions, which is the case
that matters, is unmeasured.

\item \textbf{Denomination coverage.} Results aggregate across seven
denominations, and we do not break performance out by denomination. The corpus
also retains \rupee2000 images, a denomination the Reserve Bank has since
withdrawn from circulation, which dilutes the benchmark's relevance by an
unmeasured amount.

\item \textbf{Single architecture.} We ran one configuration and report it. No
ablation isolates the contribution of the custom head, the weighted sampler, or
the individual augmentations.

\item \textbf{Latency measured off-target.} We timed inference on a desktop
Ryzen 9 7900X, not on the i5-1335U that trained the model and not on a phone
SoC, which is the deployment target the design assumes.
\end{itemize}

\section{Conclusion}

TrustNote-AI classifies Indian banknote photographs as genuine or counterfeit at
96.97\% accuracy and 0.9959 ROC-AUC on a 1121-image held-out set, in a 14.3~MB
model that trains in 36 minutes on a laptop CPU and answers in 9.1~ms. The
engineering result is solid, and the pipeline is reproducible end to end on
hardware a student already owns.

The scientific result is more interesting than the accuracy figure and points in
a less comfortable direction. Those numbers came from a network whose every
convolutional filter was fitted on ImageNet photographs of animals and household
objects, with only a half-million-parameter head trained on currency. Security
features are not in that representation. Whatever separates these two classes at
0.9959 AUC is coarse enough for generic ImageNet features to capture it, which
places print-process artifacts and capture conditions ahead of intaglio texture
as the likely explanation.

Separating those two possibilities requires the cross-source evaluation
described in Section~\ref{sec:protocol}, and until that experiment runs, the
honest claim is narrower than the headline. The model separates the classes in
this corpus. Whether it detects counterfeits remains open.

\bibliographystyle{IEEEtran}
\bibliography{references}

\end{document}
